\documentclass[a4paper,11pt]{article}
\usepackage[bottom=0.8in,top=0.8in, left=0.8in,right=0.8in]{geometry}
\pagestyle{empty}

\usepackage[T1]{fontenc}
\usepackage{newtxmath,newtxtext}
\usepackage{fontawesome5}
\usepackage{orcidlink}
\usepackage{parskip}
\usepackage{titlesec}
\usepackage{enumitem}
\usepackage{hyperref}
\usepackage[usenames,dvipsnames]{xcolor}

\definecolor{linkcolour}{rgb}{0,0.2,0.6}
\hypersetup{colorlinks,breaklinks,urlcolor={linkcolour!90}, linkcolor={linkcolour!90}}

\linespread{1.0}

\newcommand{\sectitle}[1]{
    \vspace{2ex}
    {\Large\textbf{#1}}
    \vspace{0.5ex}
    \hrule
    \vspace{1.5ex}
}

\newenvironment{cvcontent}{
    \begin{list}{}{
        \leftmargin=1.5em
        \labelwidth=0pt
        \labelsep=0pt
    }\item[]
}{
    \end{list}
}

\newsavebox{\CBox}
\def\textBF#1{\sbox\CBox{#1}\resizebox{\wd\CBox}{\ht\CBox}{\textcolor{linkcolour}{{#1}}}}

% Samsung Research Future Robotics version
\begin{document}

\noindent
\vspace{3pt}{\Huge \textbf{Seongmin Hong}} \hfill \href{mailto:smhongok@snu.ac.kr}{\faEnvelope \ smhongok@snu.ac.kr} \\
{\large Postdoctoral Researcher, INMC, Seoul National University} \hfill \href{https://smhongok.github.io}{\faGlobe \ Website} \quad \href{https://linkedin.com/in/sm-hong}{\faLinkedin \ LinkedIn} \quad \href{https://scholar.google.com/citations?user=YOUR_ID}{\faGraduationCap \ Scholar}

\vspace{0.5ex}
\hrule height 0.8pt
\vspace{1ex}

\noindent
My research lies at the intersection of \textbf{computational imaging}, \textbf{inverse problems}, and \textbf{deep generative modeling}, with a focus on robust sensing and perception for \textbf{real-world deployment}.
I specialize in developing \textbf{optimization-driven methods} that integrate physical constraints into generative architectures, producing numerically reliable outputs applicable to \textbf{embodied} and \textbf{robotic systems}.
My doctoral research on the inversion of deep generative models was recognized with the \textbf{Distinguished Ph.D. Dissertation Award}.

\vspace{2pt}

\noindent\textbf{Research Areas:}
Computer Vision \ $\cdot$ \ Machine Learning \ $\cdot$ \ Computational Imaging \ $\cdot$ \ Inverse Problems \ $\cdot$ \ Optimization \ $\cdot$ \ 3D Sensing \ $\cdot$ \ Radar Signal Processing

\vspace{2pt}

\sectitle{Education}
\begin{cvcontent}
\textbf{Ph.D.} in Electrical and Computer Engineering -- {Seoul National University}\hfill {03 2020 --- 08 2025}\\
\textit{Dissertation}: Inversion of Deep Generative Models: Structural and Optimization-based Approaches\\
\textit{Award}: \textbf{Distinguished Ph.D. Dissertation}\\
Advisor: Se Young Chun, GPA: 4.08/4.3

\textbf{B.S.} in Electrical and Computer Engineering -- {Seoul National University}  \hfill {03 2016 --- 02 2020}\\
GPA: 3.70/4.3

Math, Physics, and Astronomy -- {Gyeonggi Science High School for the Gifted} \hfill {03 2013 --- 02 2016}

\end{cvcontent}
\sectitle{Experience}
\begin{cvcontent}
\textbf{Postdoctoral Researcher (Alt. military service)} -- {INMC, SNU}\hfill {09 2025 --- 08 2026}

\end{cvcontent}
\sectitle{Honors \& Fellowships}
\begin{cvcontent}
Apple Scholars in AI/ML PhD fellowship program - Institutional Nominee \hfill 2025

Qualcomm Innovation Fellowship Korea 2024 (won 4,000,000 KRW) \hfill 12 2024

Brain Korea 21 Four Program \hfill 03 2020 --- 08 2025\\

\end{cvcontent}
\sectitle{Top-tier Conference Publications (First-authored)}
\begin{cvcontent}
\href{https://diffbmp.com/}{DiffBMP: Differentiable Rendering with Bitmap Primitives}\\
\textbf{Seongmin Hong}*, Junghun James Kim*, Daehyeop Kim, Insoo Chung, Se Young Chun (*co-first authored)\\
Conference on Computer Vision and Pattern Recognition (\textbf{CVPR}), 2026.

\href{https://smhongok.github.io/dec-inv.html}{Gradient-free Decoder Inversion in Latent Diffusion Models}\\
\textbf{Seongmin Hong}, Suh Yoon Jeon, Kyeonghyun Lee, Ernest K. Ryu, Se Young Chun\\
Neural Information Processing Systems (\textbf{NeurIPS}), 2024.

\href{https://smhongok.github.io/ada-sel.html}{Adaptive Selection of Sampling-Reconstruction in Fourier Compressed Sensing}\\
\textbf{Seongmin Hong}, Jaehyeok Bae, Jongho Lee, Se Young Chun\\
European Conference on Computer Vision (\textbf{ECCV}), 2024.

\href{https://smhongok.github.io/inv-dpm.html}{On Exact Inversion of DPM-Solvers}\\
\textbf{Seongmin Hong}, Kyeonghyun Lee, Suh Yoon Jeon, Hyewon Bae, Se Young Chun\\
Conference on Computer Vision and Pattern Recognition (\textbf{CVPR}), 2024.

\href{https://smhongok.github.io/robustness.html}{On the Robustness of Normalizing Flows for Inverse Problems in Imaging}\\
\textbf{Seongmin Hong}, Inbum Park, Se Young Chun\\
International Conference on Computer Vision (\textbf{ICCV}), 2023.

\href{https://smhongok.github.io/nubsf.html}{Neural Diffeomorphic Non-uniform B-spline Flows}\\
\textbf{Seongmin Hong}, Se Young Chun\\
Thirty-Seventh AAAI Conference on Artificial Intelligence (\textbf{AAAI}, oral), 2023.

\end{cvcontent}
\sectitle{Other Conference Publications}
\begin{cvcontent}
Triadic Dynamics Aware Diffusion Posterior Sampling for Inverse Problems: Optimizing Guidance and Stochasticity Schedules\\
Junseo Bang$^*$, Dong Ju Mun$^*$, Hoigi Seo, \textbf{Seongmin Hong}, Se Young Chun\\
International Conference on Machine Learning (\textbf{ICML}), 2026.

\end{cvcontent}
\sectitle{International Journal Publications}
\begin{cvcontent}
\href{https://www.sciencedirect.com/science/article/pii/S0167865525003083}{Robust Watermarks for Audio Diffusion Models by Quadrature Amplitude Modulation}\\
Kyoungryeol Lee$^*$, \textbf{Seongmin Hong}$^*$, Se Young Chun (*co-first authored)\\
Pattern Recognition Letters, Vol. 198, pp. 22-28, 2025.

\href{https://doi.org/10.1016/j.imavis.2025.105786}{Fidelity-Preserving Zero-shot Diffusion Models for Highly Ill-posed Inverse Problems in Lensless Imaging}\\
Haechang Lee, Dong Ju Mun, Hyunwoo Lee, Kyung Chul Lee, \textbf{Seongmin Hong}, Gwanghyun Kim, Seung Ah Lee, Se Young Chun\\
Image and Vision Computing, Volume 165, January 2026.

\href{https://ietresearch.onlinelibrary.wiley.com/doi/10.1049/rsn2.12319}{Advanced direction of arrival estimation using step-learnt iterative soft-thresholding for frequency-modulated continuous wave multiple-input multiple-output radar}\\
\textbf{Seongmin Hong}, Seong-Cheol Kim, Seongwook Lee\\
IET Radar, Sonar \& Navigation, 2023.

\href{https://ieeexplore.ieee.org/abstract/document/9552897}{Radar Signal Decomposition in Steering Vector Space for Multi-Target Classification}\\
\textbf{Seongmin Hong}, Seongwook Lee, Byeong-Ho Lee, Jinwook Kim, Yong-Hwa Kim, Seong-Cheol Kim\\
IEEE Sensors Journal, 2021.

\end{cvcontent}
\sectitle{Preprints \& Manuscripts under Review}
\begin{cvcontent}
    \textBF{DiffStamp: Aesthetically Visible Watermarking by Integrating Differentiable Renderer into Diffusion Models} \\
    \textbf{Seongmin Hong}$^*$, Kyungryeol Lee$^*$, Junghun James Kim$^*$, Se Young Chun, 2026 \\

    \href{https://dx.doi.org/10.2139/ssrn.5397023}{Principled Reframing in Motion Transfer with Video Diffusion Models} \\
    Donggyu Lee*, \textbf{Seongmin Hong}*, Se Young Chun (*co-first authored), 2025 \\
\end{cvcontent}

\sectitle{Invited Talks}
\begin{cvcontent}
\textBF{Inversion of Deep Generative Models: Structural and Optimization-based Approaches}\\
At Chung-Ang University, Seoul, Korea, Oct 22, 2025

\textBF{Inversion of Deep Generative Models: Structural and Optimization-based Approaches}\\
At ETRI, Daejeon, Korea, Oct 21, 2025

\textBF{On the Robustness of Normalizing Flows for Inverse Problems in Imaging}\\
In 2024 SIAM Conference on Imaging Science (\textbf{SIAM IS24}), Atlanta, Georgia, U.S., May 28-31, 2024

\end{cvcontent}
\sectitle{Exhibitions \& Art Projects}
\begin{cvcontent}
\href{https://yeyak.seoul.go.kr/web/reservation/selectReservView.do?rsv_svc_id=S260206141656111035}{Engineer's AI Media Art: Digital Assemblage} \hfill {02 2026}\\
Led a special workshop at Seoul Science Museum demonstrating DiffBMP technology to create digital assemblage artworks. Taught participants how to use simple recognition methods to transform everyday objects into artistic expressions, exploring the intersection of engineering, technology, and art.

\href{https://culture.snu.ac.kr/event/ipa}{IPA (Interactive Perspective Art): Ciao, Technology} \hfill {01 2026}\\
Collaborative kinetic art exhibition at SNU Powerplant. Served as the \textbf{lead for software infrastructure and typography} for the signature piece \textit{Ordered Volume}---a robotic sculpture utilizing 160 motors and spherical voxels to translate digital data into physical forms.

\href{https://tumblbug.com/arrco}{ARRCO: Multi-view 3D Crystal Point-Cloud Art} \hfill {2025}\\
Independent crowd-funded project. Developed the full computational pipeline for high-density point-cloud laser engraving to create multi-view optical illusions where distinct images emerge from different observer perspectives.

\end{cvcontent}

\sectitle{Patents}
\begin{cvcontent}
\textBF{Audio Data Management Device and Method Using A Generative AI Model}
\\Se Young Chun, \textbf{Seongmin Hong}, Kyungryeol Lee\\
Korea Patent Filed (Application No. 10-2025-0046656), 04 2025

\end{cvcontent}
\sectitle{Teaching Experience}
\begin{cvcontent}
Head TA, Topics in Signal Processing (Generative Diffusion Model, English Lecture) \hfill 2025 Autumn

Head TA, Signals and Systems (English Lecture) \hfill 2022 Spring

Head TA, Seminar in Electrical and Computer Engineering 3 \hfill 2020 Spring

TA, Advanced Electromagnetics 1 \hfill 2020 Spring
\end{cvcontent}

\sectitle{Mentoring \& Advising}
\begin{cvcontent}
\textbf{Research Mentor for 10+ Undergraduate Interns} \hfill 2022 --- Present \\
\textit{Intelligent Computational Imaging Lab (ICL), SNU}

\begin{itemize}[leftmargin=1.5em, itemsep=0.3ex]
    \item \href{https://jaehyeokbae.me/}{Jaehyeok Bae} (Now Ph. D. student at \textbf{Stanford University}) --- \textit{ECCV 2024}
    \item \href{https://inbumpark.github.io/}{Inbum Park} (Now Ph. D. student at \textbf{University of Maryland}) --- \textit{ICCV 2023}
    \item \href{https://www.linkedin.com/in/insoo-chung-07a242358/}{Insoo Chung} (Now Research Scientist at \textbf{AIRS Medical})
    \item \href{https://scholar.google.com/citations?user=guC7J3YAAAAJ&hl=ko}{Kyeonghyun Lee} (Now Ph.D. student at \textbf{Seoul National University}) --- \textit{CVPR 2024}
    \item \href{https://scholar.google.com/citations?user=Xei6FKcAAAAJ&hl=ko}{Kyungryeol Lee} (Now Ph.D. student at \textbf{Seoul National University}) --- \textit{PRL 2025}
    \item Dohyun Mah (Now M.S. student at \textbf{Seoul National University})
    \item Donggyu Lee (Undergraduate Intern) --- \textit{\href{https://dx.doi.org/10.2139/ssrn.5397023}{Preprint} 2025}
    \item \textbf{Mentored 3+ additional interns} on their early-stage research and career development.
\end{itemize}
\end{cvcontent}

\sectitle{Academic Services}
\begin{cvcontent}
\href{https://elaw.klri.re.kr/kor_service/jomunPrint.do?hseq=10572&cseq=235795}{Technical Research Personnel} (Alternative military service)
\hfill 03 2023 --- 02 2025, 09 2025 --- 08 2026

Reviewer: NeurIPS (2024, 2025), ICML (2025, \textbf{Gold 2026}), ICLR (2025, 2026), CVPR (2025), ICCV (2025), \\ AAAI (2026), AISTATS (2025, 2026), ACMMM (2025)\\ IEEE Sensors Journal, ACM TOMM, IEEE TPAMI, IEEE TMM \hfill since 2023

{\small $\cdot$ ICML 2026 Gold Reviewer: awarded to top 25\% of reviewers based on area chair ratings}

\end{cvcontent}
\sectitle{Challenges}
\begin{cvcontent}
\href{https://codalab.lisn.upsaclay.fr/competitions/17339}{AIM 2024 Depth Upsampling Challenge} - 3rd place\\
Woojae Han, Kyeonghyun Lee, \textbf{Seongmin Hong}, Se Young Chun\\

\end{cvcontent}
\sectitle{Workshop Papers}
\begin{cvcontent}
\href{https://arxiv.org/abs/2603.10445}{\textBF{Unlearning the Unpromptable: Prompt-free Instance Unlearning in Diffusion Models}}\\
Kyungryeol Lee$^*$, Kyeonghyun Lee$^*$, \textbf{Seongmin Hong}$^*$, Byung Hyun Lee, Se Young Chun\\
Machine Unlearning for Vision (MUV) Workshop in CVPR (extended abstract), 2026.

\href{https://uncertainty-cv.github.io/2024/papers/}{Uncertainty-Adaptive Selection of Sampling-Reconstruction for Fourier Compressed Sensing}\\
\textbf{Seongmin Hong}, Jaehyeok Bae, Jongho Lee, Se Young Chun\\
3rd Workshop on Uncertainty Quantification for Computer Vision in ECCV (extended abstract), 2024.

\end{cvcontent}
\sectitle{Domestic Publications}
\begin{cvcontent}
\href{https://journal-home.s3.ap-northeast-2.amazonaws.com/site/2020kics/presentation/0387.pdf}{Optimization of 16-QAM Constellations with AWGN Channel}\\
\textbf{Seongmin Hong}, Jeong-Hoon Park, Doyoung Ham, Seong-Cheol Kim\\
KICS Summer Conference (in Korean), 2020.

\textBF{Target classification based on point cloud data using automotive imaging radar}\\
Sohee Lim, Jaehoon Jung, \textbf{Seongmin Hong}, Byeong-ho Lee, Seong-Cheol Kim\\
KICS Winter Conference (in Korean), 2021.

\textBF{Measurement Data Release on 2.4 GHz ISM Band Usage}\\
Doyoung Ham, \textbf{Seongmin Hong}, Byonghyo Shim, Seong-Cheol Kim\\
KICS Winter Conference (in Korean), 2021.

\textBF{Analysis of CNN Classifier Performance for Segway Data with 77GHz Automotive Radar}\\
Byeong-ho Lee, \textbf{Seongmin Hong}, Jinwook Kim, Seongwook Lee, Byonghyo Sim, Seong-Cheol Kim\\
KICS Winter Conference (in Korean), 2021.

\end{cvcontent}
\sectitle{References}
\begin{cvcontent}
    \begin{minipage}[t]{0.5\textwidth}
        \textbf{Prof. Se Young Chun} (Advisor)\\
        Professor, ECE\\
        Seoul National University\\
        \href{mailto:sychun@snu.ac.kr}{sychun@snu.ac.kr}
    \end{minipage}
    \begin{minipage}[t]{0.5\textwidth}
        \textbf{Prof. Ernest K. Ryu}\\
        Assistant Professor, Applied Mathematics\\
        UCLA\\
        \href{mailto:eryu@math.ucla.edu}{eryu@math.ucla.edu}
    \end{minipage}
    \\
    \\
    \\
    \begin{minipage}[t]{0.5\textwidth}
        \textbf{Prof. Jongho Lee}\\
        Professor, ECE\\
        Seoul National University\\
        \href{mailto:jonghoyi@snu.ac.kr}{jonghoyi@snu.ac.kr}
    \end{minipage}
\end{cvcontent}


\end{document}
